\documentclass{article}
\usepackage{geometry}
\usepackage{concmath}
\usepackage[OT1]{fontenc}
\usepackage{graphicx}
\usepackage{amsfonts}
\usepackage{amsmath,mathtools}
\usepackage{calrsfs}
\usepackage{color}
\usepackage[dvipsnames]{xcolor}
\usepackage[many]{tcolorbox}
\usepackage{setspace}
\usepackage{tkz-tab}    
\usepackage{amssymb}          
\usepackage{multicol}  
\usepackage{esint}
\geometry{margin=0.5in}
\definecolor{green}{RGB}{52, 235, 168}
\definecolor{violet}{RGB}{165, 83, 237}
\definecolor{lemon}{RGB}{170, 237, 83}
\definecolor{red}{RGB}{240, 125, 84}
\definecolor{blue}{RGB}{105, 255, 242}

\begin{document} \noindent
\colorbox{green}{250.2}\\
\(\lambda\) valeur propre de u \(\Rightarrow P(\lambda)\) valeur propre de P(u)\\

\noindent \colorbox{green}{250.8}\\
\begin{align*}
    M \text{ inversible } &\Leftrightarrow 0 \notin Sp(M)\\
    &\Leftrightarrow 0 \text{ n'est pas racine de } \Pi_M\\
    &\Leftrightarrow \Pi_m(0) \neq 0
\end{align*}

\noindent \colorbox{green}{250.9}\\
A/ \(SpA = \{2,3\}\)\\ 
donc \((X-2)(X-3)|\Pi_A\) et \(\chi_A(A) = 0\)\\
Donc \(\Pi_A = (X-2)(X-3)\)\\
div euclidienne de \(X^n\) par \(\Pi_A\)\begin{align*}
    X^n &= Q(X-2)(X-3) + \alpha_n X + \beta_n \\
    &\begin{cases}
        2^n = 2\alpha_n + \beta_n\\
        3^n = 3\alpha_n + \beta_n
    \end{cases}\\
    A^n &= Q(A)\Pi_A(A) + \alpha_n A + \beta_n I_2\\
    &= \alpha_n A + \beta I_2
\end{align*}
B/ \(\chi_B = (X-2)^2\) 
\(B-2I_2 \neq 0\) donc \(\Pi_b = (X-2)^2\)\begin{align*}
    X^n = Q(X-2)^2 + \alpha_n X + \beta_n \\
    nX^{n-1} = (X-2)qqch + \alpha_n
\end{align*}



\end{document}